%% The following is a directive for TeXShop to indicate the main file
%%!TEX root = diss.tex

%=======================================
\chapter{OKM in Highly symmetric 2D models}
%=======================================
We'll continue further with the OKM for highly-symmetric models. There are only a few such models. If $x_k$ is in $S$ then so is $\ol{x}_k$. If one of the diagonal terms is in $S$ then so are the others. To have a group of order 4 we need to have both an $x_1$ and $x_2$ variable in $S$ so we have several non-isomorphic models. One with cartesian steps, one with diagonal steps, one with diagonal step and $x_1$,$\ol{x}_1$ and finally the model with all steps. The notation will be the same as in previous chapters. We derived the following equations
\begin{align}
	x_1 = \sigma_1:\quad\Delta_2(K2^\ast Q_1^0) &=\frac{1}{K_1} \left(\Delta_2(K_2^\ast K_{12})\Qf{0,0} - \alpha(K_2^\ast-K_2)\right)\,,\\
	x_2 = \sigma_2:\quad\Delta_1(K1^\bullet Q_2^0) &=\frac{1}{K_2} \left(\Delta_1(K_1^\bullet K_{12})\Qf{0,0} - \alpha(K_1^\bullet-K_1)\right)\,.
\end{align}
It is enough to only consider the analysis for one of the quations for now. Only the third model requires looking at both equations. We already know from the previous chapter that $\mu_2 = S_2^{\um 1}/S_2^1 = 1$ and 
\begin{align}
	\alpha(K_2^\ast -K_2)  &= \alpha t(\ol{x}_2 - x_2)S_2^{\um 1}(\sigma_1)\,,\\
	\Delta_2(K_2^\ast K_{12}) &= c_2t(\ol{x}_2-x_2)(c_1S_2^{\um 1}(\sigma_1) - \ol{\sigma}_1\Sf{\um 1,\um 1})\,.
\end{align}
To compute $\Delta_2(K2^\ast Q_1^0)$ we need to compute $K_2(\sigma_1,x_2)$ and $K_2^\ast(\sigma_1,x_2)$ which are given by
\begin{align}
	K_2 = c_2 -t\ol{x}_2S_2^{\um 1}(\sigma_1) \com K_2^\ast = c_2 -t x_2S_2^{\um 1}(\sigma_1)\,.
\end{align}
Therefore we will focus on computing $S_2^{\um 1}(\sigma)$ first.
\begin{align}
\begin{split}
	S_2^{\um 1}(\sigma_1) &= \Sf{0,1}+ \Sf{1,1}(\ol{\sigma}_1+\sigma_1)\\
	&= \Sf{0,1} + \Sf{1,1}\frac{1-tS_1^0}{t S_1^1}\\
	&= \frac{1}{tS_1^1}\left(\Sf{1,1}+t\left(\Sf{0,1}S_1^1-\Sf{1,1}S_1^0\right)\right)\\
	&= \frac{\Sf{1,1}+t\Sf{0,1}\Sf{1,0}}{t\left(\Sf{1,0}+\Sf{1,1}(\ol{x}_2+x_2)\right)}\,.
\end{split}
\end{align}
Therefore the function $S_2^{\um 1}(\sigma_1)$ is a rational function of $x_2$.
\begin{table}[H]
\caption{Function $S_2^{- 1}(\sigma_1)$ for highly symmetric models.}
\label{tbl:Symmetric}
\centering
\begin{tabular}{ccccc}
\toprule
	Model & $\Sf{0,1}$ & $\Sf{1,0}$ & $\Sf{1,1}$ & $S_2^{\um 1}(\sigma_1)$\\
\midrule
	\cartmodel & 1 & 1 & 0 & 1 \\
	\diagmodel & 0 & 0 & 1 & $\ol{t}/(\ol{x}_2+x_2)$\\
	\modelgraph 11100111 & 1 & 0 & 1 & $\ol{t}/(\ol{x}_2+x_2)$\\
	\modelgraph 10111101 & 0 & 1 & 1 & $\ol{t}/(\ol{x}_2+1+x_2)$\\
	\fullmodel & 1 & 1 & 1 & $(1+\ol{t})/(\ol{x}_2+1+x_2)$\\
\bottomrule
\end{tabular}
\end{table}
\begin{lemma} For highly symmetric models with  $e = \Sf{1,1}, f = \Sf{1,0}, g = \Sf{0,1}$ we have
\begin{align}
	\frac{\alpha(K_2^\ast -K_2)}{K_1}  &= \frac{\alpha (\ol{x}_2 - x_2)(e+tfg)}{S_1^1K_1}\,,\\
	\frac{\Delta_2(K_2^\ast K_{12})}{K_1} &= \frac{c_2(\ol{x}_2-x_2)(c_1(e+tfg)- et \ol{\sigma}_1S_1^1)}{S_1^1(c_1-t\ol{\sigma}_1S_1^1)}\,.
\end{align}
\end{lemma}
%=======================================
\section{Cartesian model: steps N, S, E, W}
%=======================================
In this case $S_2^{\um 1}(\sigma_1) = 1$ so a lot of the calculations will be simpler and this is a good place to start further work. Here $K_{12} = c_1c_2$ so it behaves like $\alpha$ and we have
\begin{align}
	\Delta_2(K2^\ast Q_1^0) + \frac{\alpha\beta(K_2^\ast-K_2)}{K_1}=0\,,
\end{align}
where
\begin{align}
	\beta =  1-b_1b_2 Q_{1,2}^{0,0}
\end{align}
and
\begin{align}
	\frac{\alpha\beta(K_2^\ast -K_2)}{K_1}  &= \frac{\alpha\beta t (\ol{x}_2 - x_2)}{K_1} = \alpha\beta f(c_1,x_2)\,,
\end{align}
since $K_1$ doesn't depend on $c_2$ or $x_1$. The equation finally becomes
\begin{align}
	(c_2-tx_2)Q_1^0(x_2) - (c_2-t\ol{x}_2)Q_1^0(\ol{x}_2) + \alpha\beta f = 0\,.
\end{align}
We perform the extraction $[x_2^1]$ to get
\begin{align}
	c_2Q_{1,2}^{0,1}-t Q_{1,2}^{0,0} + \alpha\beta f_2^1 = 0
\end{align}
and the extraction $[b_2^0]$ to get
\begin{align}
	- t \wc{Q}_{1,2}^{0,0} +  \ol{a}_1f_2^1\,.
\end{align}
Using the second equation within the first one and simplifying yields
\begin{align}
	b_2 Q_{1,2}^{0,1} &= t(a_2 Q_{1,2}^{0,0} - \beta \wc{Q}_{1,2}^{0,0})\,,\\
	b_1 Q_{1,2}^{1,0} &= t(a_1 Q_{1,2}^{0,0} - \beta \wh{Q}_{1,2}^{0,0})
\end{align}
Here the second equation is obtained via the same method but with $\bullet$ instead of $\ast$. Using the same method we can also obtain
\begin{align}
	b_2 \Qf{1,1} = t(a_2 \Qf{1,0}-\beta \wc{Q}_{1,2}^{1,0})\,.
\end{align}
Multiplying by $b_1$ and using the above relations leads to
\begin{align}
\begin{split}
	b_1b_2 \Qf{1,1} &= t(a_2b_1\Qf{1,0}-\beta b_1 \Qc{1,0})\\
			      &=t a_2t(a_1\Qf{0,0}-\beta \Qh{0,0})-t\beta t(a_1\Qc{0,0}-\beta\Qch{0,0})\\	
			      &=t^2(a_1a_2\Qf{0,0}-a_2\beta\Qh{0,0}-a_1\beta\Qc{0,0}+\beta^2\Qch{0,0})\,.
\end{split}
\end{align}
Let
\begin{align}
	G = \Qf{0,0}, \quad\wh{G} = G_1, \quad\wc{G} = G_2\,.
\end{align}
From before we have
\begin{align}
	\Qf{0,0} = 1 + t a_1a_2(\Qf{1,0}+\Qf{0,1})\,.
\end{align}
Multiplying by $b_1b_2$ and using the above relations leads to
\begin{align}
\begin{split}
b_1b_2 G &= b_1b_2 + ta_1a_2(b_2b_1\Qf{1,0}+b_1b_2\Qf{0,1})\\
                 &=b_1b_2 + ta_1a_2(b_2t(a_1G-\beta G_1)+b_1t(a_2 G- \beta G_2)\\
                 &=b_1b_2 + t^2a_1a_2(b_2a_1G+b_1a_2G-b_2\beta G_1-b_1\beta G_2)\,.
\end{split}
\end{align}
Using the relation for $\beta$ and simplifying leads to
\begin{align}
	G = \frac{b_1b-2 - t^2 a_1a_2(b_2G_1+b_1G_2)}{b_1b_2 - t^2a_1a_2 (a_2b_1+a_1b_2+b_1b_2(b_2G_1+b_1G_2))}\,.
\end{align}
If we let $F = (b_2G_1+b_1G_2)$ and perform the long division with respect to $F$ we find that
\begin{align}
	\Qf{0,0} = \frac{1}{b_1b_2}\left(1-\frac{b_1b_2(1-b_1b_2)-t^2a_1a_2(a_1b_2+a_2b_1)}{b_1b_2+t^2a_1a_2(a_1b_2+a_2b_1+b_1b_2 F)}\right)\,.
\end{align}
Using arguments similar to those in the later section on diagonal steps we can find out more about the various $Q$'s.
\begin{align}
	b_2\Qf{m,n+1} &= t(a_2\Qf{m,n}-\beta \Qc{m,n})\,,\\
	b_1\Qf{m+1,n} &= t(a_1\Qf{m,n}-\beta \Qh{m,n})\,.
\end{align}
Summing the equations multiplied by their respective powers of $x_1$ and $x_2$ we find the equations
\begin{align}
	b_2 Q_2^{n+1} &= t(a_2Q_2^n-\beta \wc{Q}_2^n)\,,\\
	b_1 Q_1^{m+1} &= t(a_1Q_1^m-\beta \wh{Q}_1^m)\,.
\end{align}
Now we can sum with respect to $x_2$ and $x_1$ to find
\begin{align}
	b_2 (Q-Q_2^0) &= x_2t(a_2Q - \beta \wc{Q})\,,\\
	b_1 (Q-Q_1^0) &= x_1t(a_1Q - \beta \wh{Q})\,.
\end{align}
We can solve these for $Q_2^0$ and $Q_1^0$ as follows
\begin{align}
	b_2Q_2^0 &= (b_2-a_2x_2t)Q + x_2t\beta \wc{Q}\,,\\
	b_1Q_1^0 &= (b_1-a_1x_1t)Q + x_1t\beta \wh{Q}\,.
\end{align}
We substitute these two expressions into the original functional equation and after some algebra we find that
\begin{align}
	Q = \frac{(1-b_1b_2 G)(\alpha + tx_1\ol{b_1}K_1\wh{Q}+tx_2\ol{b_2}K_2 \wc{Q})}{K-(1-\ol{c_1}x_1t)K_1-(1-\ol{c_2}x_2t)K_2}
\end{align}
and after recognizing some of the quantities as modified $K_1$ and $K_2$'s we finally stop at the solution for $Q$
\begin{align}
	Q = \frac{(1-b_1b_2 \Qf{0,0})(\alpha - \ol{b}_1 K_1 \wh{K}_1^\bullet \wh{Q}-\ol{b}_2K_2\wc{K}_2^\ast \wc{Q})}{K - \ol{c}_1 K_1 K_1^\bullet - \ol{c}_2 K_2 K_2^\ast}\,.
\end{align}
This rather complicated expression shows that the generating function$Q$ is obtained via copies of $Q$ which have some of their marking variables set to 1.
%=======================================
\section{Diagonal model: steps NE, NW, SE, SW}
%=======================================
The next case in terms of difficulty is when we have diagonal steps. This model is very similar to the one covered in \cite{wong2014confining} where they have two interacting walks confined in a semi-infinite strip. The models are different because there we have an extra term in the functional equations. The difficulty in the current case perhaps arrises in that $S_2^{\um 1}(\sigma_1)$ is now more complicated which is due to the diagonal terms. Also $\Sf{\um 1, \um 1} \neq 1$ so $K_{12}$ is not trivial. Here
\begin{align}
	\frac{\alpha(K_2^\ast -K_2)}{K_1}  &= \frac{\alpha (\ol{x}_2 - x_2)}{K_1(\ol{x}_2+x_2)}\,,\\
	\frac{\Delta_2(K_2^\ast K_{12})}{K_1} &= \frac{c_2(\ol{x}_2-x_2)(c_1- t \ol{\sigma}_1S_1^1)}{S_1^1(c_1-t\ol{\sigma}_1S_1^1)} = \frac{c_2(\ol{x}_2-x_2)}{\ol{x}_2+x_2}\,.
\end{align}
The functional equation becomes
\begin{align}
   K_2^\ast Q_1 - K_2 Q_1^\ast -\frac{c_2(\ol{x}_2-x_2)}{\ol{x}_2+x_2}Q_{1,2}+\alpha\frac{\ol{x}_2-x_2}{(\ol{x}_2+x_2)K_1}= 0\,.
\end{align}
Surprisingly enough, the only term that depends on $\sigma_1$ remaining in this expression is $K_1$ and this will be of no consequence to the rest of the analysis. We just need to take care of the remaining terms. We'll do this by multiplying the equation by $\ol{x}_2+x_2$. We have
\begin{align}
	R^\ast \exx= (\ol{x}_2+x_2) K_2^\ast &= c_2(\ol{x}_2+x_2)-x_2 = c_2\ol{x}_2- \ol{a}_2x_2\,,\\
	R = (\ol{x}_2+x_2) K_2 &= c_2(\ol{x}_2+x_2)-\ol{x}_2 = c_2x_2- \ol{a}_2\ol{x}_2\,.
\end{align}
There are a couple of relations here
\begin{align}
	R^\ast - R &= v(K_2^\ast-K_2) = \ol{x}_2-x_2\\
	\ol{x}_2R^\ast - x_2 R &= \ol{x}_2c_2v-1- x_2 c_2v+1 \\
				      & = c_2v(\ol{x}_2-x_2)
\end{align}
The equation becomes
\begin{align}
(c_2\ol{x}_2- \ol{a}_2x_2)Q_1^0(x_2)+ ( \ol{a}_2\ol{x}_2-c_2x_2) Q_1^0(\ol{x}_2)+(\alpha \ol{K}_1 - c_2Q_{1,2}^{0,0})(\ol{x}_2-x_2)= 0
\end{align}
and it's ready for extraction. Let $\delta_0 =  \ol{a}_1(\ol{x}_2-x_2)\ol{K}_1$. We re-write the equation as
\begin{align}\label{zero-order}
    R^\ast Q_1^0 -R Q_1^{0,\ast} - c_2(\ol{x}_2-x_2)\Qf{0,0} + \ol{a}_2\delta_0 = 0\,.
\end{align}
We'll use the master equation
\begin{align}
(1-tS)Q - (c_1-t\ol{x}_1S_1^{\um 1}) Q_1^0- (c_2-t\ol{x}_2S_2^{\um 1}) Q_2^0+ (c_1c_2-t\ol{x}_1\ol{x}_2S_{1,2}^{\um 1,\um 1})Q_{1,2}^{0,0}  = \ol{a}_1\ol{a}_2
\end{align}
to further our work here. For diagonal steps it is
\begin{align}
 (1-tS)Q - (c_1-tv\ol{x}_1) Q_1^0- (c_2-t\ol{x}_2S_2^{\um 1}) Q_2^0+ (c_1c_2-t\ol{x}_1\ol{x}_2)Q_{1,2}^{0,0}  = \ol{a}_1\ol{a}_2
\end{align}
extracting the $[x_1^0]$ coefficient yields
\begin{align}
   Q_1^0 =  \ol{a}_2+c_2\Qf{0,0} + a_1tvQ_1^1\,.
\end{align}
We can extract higher coefficients too. Let $k\geq 1$. Then the $[x_2^k]$ coefficient is
\begin{align}\begin{split}
	0 &= Q_1^k-tS_{1}^{\um 1}Q_1^{k+1}-tS_1^{0}Q_1^k-tS_1^1Q_1^{k-1}-c_2\Qf{k,0} \\ 
	   &\quad+ t\ol{x}_2(\Sf{\um 1,\um 1}\Qf{k+1,0}+\Sf{0,\um 1}\Qf{k,0}+\Sf{1,\um 1}\Qf{k-1,0})\\
	   &=Q_1^k-tvQ_1^{k+1}-tvQ_1^{k-1}-c_2\Qf{k,0} + t\ol{x}_2(\Qf{k+1,0}+\Qf{k-1,0})\,.
\end{split}\end{align}
We can re-write the last one in a form more convenient for the next few steps
\begin{align}\label{formula}
	Q_1^k = tv(Q_1^{k+1}+Q_1^{k-1}) + c_2 \Qf{k,0} - t\ol{x}_2(\Qf{k+1,0}+\Qf{k-1,0})\,.
\end{align}
We'll use induction to prove a nice formula but for that we'll first need to establish another base case because of this splitting of cases for the coefficient extraction. We have
\begin{align}
	Q_1^0 &=  \ol{a}_2+c_2\Qf{0,0} + a_1tvQ_1^1\,,\\
	Q_1^{0,\ast} \exx&=  \ol{a}_2+c_2\Qf{0,0} + a_1tvQ_1^{1,\ast}\,.
\end{align}
Combining them with equation (\ref{zero-order}) gives
\begin{align}
\begin{split}
        0 &= R^\ast Q_1^0 -R Q_1^{0,\ast} - c_2(\ol{x}_2-x_2)\Qf{0,0} + \ol{a}_2\delta_0 \\
	&= R^\ast (  \ol{a}_2+c_2\Qf{0,0} + a_1tvQ_1^1) -R ( \ol{a}_2+c_2\Qf{0,0} + a_1tvQ_1^{1,\ast}) \\
	&\quad- c_2(\ol{x}_2-x_2)\Qf{0,0} + \ol{a}_2\delta_0\\
	&=a_1tv(R^\ast Q_1^1-R Q_1^{1,\ast}) + (R^\ast-R)( \ol{a}_2+c_2\Qf{0,0}) \\
	&\quad- c_2(\ol{x}_2-x_2)\Qf{0,0} + \ol{a}_2\delta_0\\
	& =a_1tv(R^\ast Q_1^1-R Q_1^{1,\ast})+ \ol{a}_2(\ol{x}_2-x_2)+ \ol{a}_2\delta_0\\
	& =a_1tv(R^\ast Q_1^1-R Q_1^{1,\ast})+ \ol{a}_2((\ol{x}_2-x_2)+\delta_0)\,.
\end{split}
\end{align}
So we can write the equation
\begin{align}
        R^\ast Q_1^1-R Q_1^{1,\ast}+ \ol{a}_2\delta_1 = 0
\end{align}
where
\begin{align}
   \delta_1 = \frac{(\ol{x}_2-x_2)+\delta_0}{a_1tv}\,.
\end{align}
This equation together with equation (\ref{zero-order}) form the base cases for the following lemma
\begin{lemma} Let $m\ge 0$ be an integer. Then
\begin{align}\label{coeff-eqn}
	R^\ast Q_1^m -R Q_1^{m,\ast} - c_2(\ol{x}_2-x_2)\Qf{m,0} + \ol{a}_2\delta_m = 0
\end{align}
where $\delta_m$ is not a function of $a_2$.
\end{lemma}
\begin{proof}We'll show this by induction. We have the base cases $m = 0$ and $m=1$. The latter one fits the pattern because $\Qf{1,0} = 0$ since we cannot have any walks end on coordinates $(1,0)$. Now we assume that the formula is true for all cases up to and including $m$. Then we use formula (\ref{formula})
\begin{align}
\begin{split}
	& R^\ast Q_1^m -R Q_1^{m,\ast}\\
	=&tv(R^\ast Q_1^{m+1}-R Q_1^{m+1,\ast}) + tv(R^\ast Q_1^{m-1}-R Q_1^{m-1,\ast})\\
	   & + c_2(R^\ast-R)\Qf{m,0}-t(\ol{x}_2R^\ast-x_2 R)(\Qf{m+1,0}+\Qf{m-1,0})\\
	=& tv(R^\ast Q_1^{m+1}-R Q_1^{m+1,\ast}) + tv(c_2(\ol{x}_2-x_2)\Qf{m-1,0}- \ol{a}_2\delta_{m-1})\\
	   & + c_2(\ol{x}_2-x_2)\Qf{m,0}-c_2tv(\ol{x}_2-x_2)(\Qf{m+1,0}+\Qf{m-1,0})\\
	=& tv(R^\ast Q_1^{m+1}-R Q_1^{m+1,\ast})+ c_2(\ol{x}_2-x_2)\Qf{m,0}\\
	   &-c_2tv(\ol{x}_2-x_2)\Qf{m+1,0}) -  \ol{a}_2tv\delta_{m-1})\,.
\end{split}
\end{align}
Combining this with equation (\ref{coeff-eqn}) leads to
\begin{align}
	0 &= tv(R^\ast Q_1^{m+1}-R Q_1^{m+1,\ast})-c_2tv(\ol{x}_2-x_2)\Qf{m+1,0} -  \ol{a}_2tv\delta_{m-1} + \ol{a}_2\delta_m\,, \\
	0 & = R^\ast Q_1^{m+1}-R Q_1^{m+1,\ast}-c_2(\ol{x}_2-x_2)\Qf{m+1,0}+  \ol{a}_2\delta_{m+1}
\end{align}
where
\begin{align}
  \delta_{m+1} =\frac{\delta_m}{tv}-\delta_{m-1}\,.
\end{align}
This function $\delta_{m+1}$ is also independent from $a_2$ and this completes the proof!
\end{proof}
There is an equivalent statement for $Q_2^n$. We can expand $R^\ast$ and $R$ to a form we can use
\begin{align}\begin{split}
	0 &=(c_2\ol{x}_2- \ol{a}_2x_2)Q_1^m(x_2) -( c_2x_2- \ol{a}_2\ol{x}_2)Q_1^m(\ol{x}_2) \\
	&\quad- c_2(\ol{x}_2-x_2)\Qf{m,0} + \ol{a}_2\delta_m \,.
\end{split}\end{align}
We start by extracting the $[x_2^1]$ coefficient
\begin{align}\begin{split}
	& c_2\Qf{m,2}- \ol{a}_2\Qf{m,0}-c_2\Qf{m,0} +c_2\Qf{m,0}+ [x_2^1] \ol{a}_2\delta_m = 0\\
	& c_2\Qf{m,2}- \ol{a}_2\Qf{m,0}+ [x_2^1] \ol{a}_2\delta_m = 0\,.
\end{split}\end{align}
We can also extract the $[x_2^{n+1}]$ coefficient for $n\geq 2$
\begin{align}
 	c_2 Q_{1,2}^{m,n+2}- \ol{a}_2Q_{1,2}^{m,n}+[x_2^{n+1}] \ol{a}_2\delta_m= 0\,.
\end{align}
This formula also contains the $n = 0$ case so we treat them together. We can set $a_2 = 1$ to get
\begin{align}
	-\wc{Q}_{1,2}^{m,n}+[x_2^{n+1}]\delta_m = 0\,.
\end{align}
Substituting this back into the above equation gives
\begin{align}
	b_2 Q_{1,2}^{m,n+2} = Q_{1,2}^{m,n} - \wc{Q}_{1,2}^{m,n}\,.
\end{align}
An equivalent formula can be derived if the process is followed through the use of $x_1$ instead of $x_2$ in the analysis.
\begin{corollary} For integers $m,n \ge 0$ we have the following relations
\begin{align}
	b_1 \Qf{m+2,n} &= Q_{1,2}^{m,n}-\wh{Q}_{1,2}^{m,n}\,,\\
	b_2 \Qf{m,n+2} &= Q_{1,2}^{m,n}-\wc{Q}_{1,2}^{m,n}\,.
\end{align}
\end{corollary}

We can sum up those equations to obtain functional equations
\begin{corollary} For integers $m \ge 0$ we have the following relations
\begin{align}
	Q_1^m - \wh{Q}_1^m &= b_1 Q_1^{m+2}\,,\\
	\wc{Q}_1^m - \wc{\wh{Q}}_1^m &= b_1 \wc{Q}_1^{m+2}\,,\\
	x_1^2(Q_2^n-\wh{Q}_2^n) &= b_1 (Q_2^n - Q_{1,2}^{1,n}x_1 - Q_{1,2}^{0,n})\,,\\
	x_1^2(\wc{Q}_2^n-\wc{\wh{Q}}_2^n) &= b_1 (\wc{Q}_2^n - \wc{Q}_{1,2}^{1,n}x_1 - \wc{Q}_{1,2}^{0,n})\,,\\
	x_1^2(Q-\wh{Q}) &= b_1(Q-x_1Q_1^1-Q_1^0)\,,\\
	x_1^2(\wc{Q}-\wc{\wh{Q}}) &= b_1(\wc{Q}-x_1\wc{Q}_1^1-\wc{Q}_1^0)\,.
\end{align}
\end{corollary}

\begin{corollary} For integers $n \ge 0$ we have the following relations
\begin{align}
	Q_2^n-\wc{Q}_2^n &= b_2 Q_2^{n+2}\,,\\
	\wh{Q}_2^n-\wc{\wh{Q}}_2^n &= b_2 \wh{Q}_2^{n+2}\,,\\
	x_2^2(Q_1^m-\wc{Q}_1^m) &=  b_2 (Q_1^m - Q_{1,2}^{m,1}x_1 - Q_{1,2}^{m,0})\,,\\
	x_2^2(\wh{Q}_1^m-\wc{\wh{Q}}_1^m) &=  b_2 (\wh{Q}_1^m - \wh{Q}_{1,2}^{m,1}x_1 - \wh{Q}_{1,2}^{m,0})\,,\\
	x_2^2(Q-\wc{Q}) &= b_2(Q-x_2Q_2^1-Q_2^0)\,,\\
	x_2^2(\wh{Q}-\wc{\wh{Q}}) &= b_2(\wh{Q}-x_2\wh{Q}_2^1-\wh{Q}_2^0)\,.
\end{align}
\end{corollary}
Every other equation is obtained from the previous one by either setting $a_1 = 1$ or $a_2 = 1$. Here we can hope to find some telescoping sum type of thing in order to say more about $Q_{1,2}^{0,0}$

Now we can move onto the main result about the diagonal walks and it is that the function for labelled walks ending at $(0,0)$ is connected to the functions for semi-labelled and unlabelled walks endind at $(0,0)$. We can use the master equation to extract the $[x_1^0x_2^0]$ and $[x_1^1x_2^1]$ coefficients yielding the following two relations
\begin{align}
	\Qf{0,0} &= 1 + ta_1a_2 \Qf{1,1}\,,\\
	\Qf{1,1} &= t(\Qf{0,0}+\Qf{2,0}+\Qf{0,2}+\Qf{2,2})\,.
\end{align}
We also have
\begin{align}
\begin{split}
	b_1b_2\Qf{2,2} &= b_1(\Qf{2,0}-\wc{Q}_{1,2}^{2,0})\,,\\
	&= b_1 \Qf{2,0} - b_1\wc{Q}_{1,2}^{2,0}\,,\\
	&= \Qf{0,0}-\wh{Q}_{1,2}^{0,0}-\wc{Q}_{1,2}^{0,0}+\wc{\wh{Q}}_{1,2}^{0,0}\,,
\end{split}
\end{align}
together with
\begin{align}
	b_1b_2\Qf{2,0} &= b_2 \Qf{0,0} - b_2 \wh{Q}_{1,2}^{0,0}\,,\\
	b_1b_2\Qf{0,2} &= b_1 \Qf{0,0} - b_1 \wc{Q}_{1,2}^{0,0}\,.
\end{align}
We we have
\begin{align}
\begin{split}
b_1b_2\Qf{1,1} &= t(b_1b_2\Qf{0,0}+b_1b_2\Qf{2,0}+b_1b_2\Qf{0,2}+b_1b_2\Qf{2,2})\\
	&= t(b_1b_2\Qf{0,0}+ b_2 \Qf{0,0} - b_2 \wh{Q}_{1,2}^{0,0}+ b_1 \Qf{0,0} \\
	&\quad- b_1 \wc{Q}_{1,2}^{0,0}+\Qf{0,0}-\wh{Q}_{1,2}^{0,0}-\wc{Q}_{1,2}^{0,0}+\wc{\wh{Q}}_{1,2}^{0,0})\\
	&= t((b_1b_2+b_1+b_2+1)\Qf{0,0}-(b_2+1)\wh{Q}_{1,2}^{0,0}\\
	&\quad-(b_1+1)\wc{Q}_{1,2}^{0,0}+\wc{\wh{Q}}_{1,2}^{0,0})\\
	&=t(a_1a_2\Qf{0,0}-a_2\wh{Q}_{1,2}^{0,0}-a_1\wc{Q}_{1,2}^{0,0}+\wc{\wh{Q}}_{1,2}^{0,0})\,,
\end{split}
\end{align}
We can finally write
\begin{align}
\begin{split}
	b_1b_2\Qf{0,0} &= b_1b_2 + ta_1a_2b_1b_2\Qf{1,1}\\
		& =  b_1b_2 + t^2a_1a_2(a_1a_2\Qf{0,0}-a_2\wh{Q}_{1,2}^{0,0}-a_1\wc{Q}_{1,2}^{0,0}+\wc{\wh{Q}}_{1,2}^{0,0})\,.
\end{split}
\end{align}
Let $G(b_1,b_2)$ stand for $Q(0,0,b_1,b_2)$ and $G_k$ has $b_k = 0$. Then we can write it all as
\begin{align}
	b_1b_2 G = b_1b_2 + t^2a_1a_2(a_1a_2G-a_2G_1-a_1G_2+G_{1,2})
\end{align}
or written as a linear combination
\begin{align}
	(b_1b_2-t^2a_1^2a_2^2)G+t^2a_1a_2^2G_1 + t^2a_1^2a_2G_2 - t^2a_1a_2G_{1,2} = b_1b_2\,.
\end{align}
We can divide out by $a_1a_2$ and obtain the equation
\begin{align}
	(c_1c_2-a_1a_2t^2)G + a_2t^2G_1+a_1t^2G_2-t^2G_{1,2}=c_1c_2\,,
\end{align}
from which we can read off
\begin{align}
	f(a_1,a_2,t) = c_1c_2-a_1a_2t^2
\end{align}
and
\begin{align}
	G = \frac{(c_1c_2-t^2(a_2G_1+a_1G_2-G_{1,2})}{f(a_1,a_2,t)}
\end{align}
where $B(fG) = c_1c_2$. Also, we can solve for $Q_k^j$ and other terms, in terms of various $[b_k^l]$ extractions, so overall it should lead to a relationship of the form $B(FQ) = 0$.

%=======================================
\section{The full model}
%=======================================
We skip right to the final model - the one with all the possible small step. Here we have the most complicated functions. Here
\begin{align}
	\frac{\alpha(K_2^\ast -K_2)}{K_1}  &= \frac{\alpha (\ol{x}_2 - x_2)(1+t)}{S_1^1K_1}\,,\\
	\frac{\Delta_2(K_2^\ast K_{12})}{K_1} &= \frac{c_2(\ol{x}_2-x_2)(c_1(1+t)- t \ol{\sigma}_1S_1^1)}{S_1^1(c_1-t\ol{\sigma}_1S_1^1)}\\
	&=\frac{c_2(\ol{x}_2-x_2))}{S_1^1}\left(1-\frac{t}{c_1-t\ol{\sigma}_1S_1^1}\right)\,.
\end{align}
We'll need to do some more work to find a representation for that last fraction. For that we'll need to $\sigma_1$ root which is
\begin{align}
\begin{split}
	\sigma_1 &= \frac{(1-tS_1^0)-\sqrt{(1-tS_1^0)^2-4t^2(S_1^1)^2}}{2tS_1^1} \\
		&= \frac{u-\sqrt{u^2-4w^2}}{2w}\\
\end{split}\\
	&\ol{\sigma}_1 = \frac{u+\sqrt{u^2-4w^2}}{2w}
\end{align}
Then we have
\begin{align}
\begin{split}
	\frac{t}{c_1-\ol{\sigma}_1w} &= \frac{2t}{2c_1-u-\sqrt{u^2-4w^2})} \\
	&= \frac{2t(2c_1-u+\sqrt{u^2-4w^2})}{(2c_1-u)^2-(u^2-4w^2)}\\
	&=  \frac{t(2c_1-u+\sqrt{u^2-4w^2})}{2(c_1^2-c_1u+w^2)}
\end{split}
\end{align}
Let $v = \ol{x}_2+x_2$. Then the denominator has
\begin{align}
\begin{split}
	c_1^2-c_1u+w^2 &= c_1^2-c_1(1-tS_1^0)+t^2(S_1^1)^2\\&=\ol{a}_1^2b_1 - \ol{a}_1b_1vt+(1+v)^2 t^2
\end{split}
\end{align}
and the numerator has
\begin{align}
\begin{split}
	u^2-4w^2 &= (1-tv)^2-4t^2(1+v)^2 \\&= 1-2tv-(3v+2)(v+2)t^2 \\&= (1-(3v+2)t)(1+(v+2)t)\,.
\end{split}
\end{align}
Thus the square root has the expansion
\begin{align}
\begin{split}
	\sqrt{u^2-4w^2} = 1 - vt-2(1+v)^2 t^2-2v(1+v)^2t^3+\cdots
\end{split}
\end{align}
We can give a series representation of the square root via the series
\begin{align}
	\sqrt{1-z} = \sum_{n=0}^\infty \frac{(2n)!}{(1-2n)(n!)^2 4^n}z^n\,.
\end{align}
That is,
\begin{align}
\begin{split}
	&\sqrt{(1+(v+2)t)(1-(3v+2)t)} \\=& \left(\sum_{n=0}^\infty \frac{(-1)^n(2n)!(v+2)^n t^n}{(1-2n)(n!)^2 4^n}\right)\left(\sum_{n=0}^\infty \frac{(2n)!(3v+2)^nt^n}{(1-2n)(n!)^2 4^n}\right)\\
	=& \sum_{n=0}^\infty \frac{t^n}{4^n}\sum_{k=0}^n \frac{(-1)^k(2k)!(2n-2k)!}{(1-2k)(1-2n+2k)k!^2 (n-k)!^2}(v+2)^k(3v+2)^{n-k}\,,
\end{split}
\end{align}
But we're not sure if a closed form exists for these coefficients. If we want to extract $x_2$ coefficients from this expression we wouldn't have an easy time. Now $K_2$ and $K_2^\ast$ will have $1+v$ in their denominators so we can multiply the whole functional equation by $1+v$ to clear out these denominator. We have
\begin{align}
	(1+v)K_2 &= c_2(1+v)-\ol{x}_2(1+t) = (c_2-1-t)\ol{x}_2+c_2+c_2x_2\,,\\
	(1+v)K_2^\ast &= c_2(1+v)-x_2(1+t) = c_2\ol{x}_2+c_2+(c_2-1-t)x_2\,,
\end{align}
and
\begin{align}
	(1+v)\frac{\alpha(K_2^\ast -K_2)}{K_1} = &\frac{\alpha (\ol{x}_2 - x_2)(1+t)}{K_1} = \alpha g(a_1,x_2,t)\,,
\end{align}
\begin{align}
\begin{split}
	(1+v)\frac{\Delta_2(K_2^\ast K_{12})}{K_1}=&c_2(\ol{x}_2-x_2)\left(1-t\frac{c_1-1+tv+\sqrt{(1+(v+2)t)(1-(3v+2)t)}}{c_1^2-c_1+c_1tv+(1+v)^2t^2}\right)\,,
\end{split}
\end{align}
and let the latter be $c_2(\ol{x}_2-x_2)f(a_1,x_2,t)$. Then if we let $A = c_2-(1+t)$ we can write the funtional equation as
\begin{align}
	(c_2\ol{x}_2+c_2+Ax_2)Q_1^0(x_2) - (A\ol{x}_2+c_2+c_2 x_2)Q_1^0(\ol{x}_2) = c_2(\ol{x}_2-x_2)f-\alpha g\,.
\end{align}
As before, we extract the $[x_2]$ coefficient from both sides to obtain
\begin{align}
	 c_2\Qf{0,2}+c_2\Qf{0,1}-(1+t)\Qf{0,0} = -\alpha g_2^1 + c_2(f_2^2-f_2^0)\Qf{0,0}\,,
\end{align}
and then set $a_2 = 1$ on both sides to find that
\begin{align}
	-(1+t)\Qc{0,0} = -\ol{a}_1 g_2^1\,.
\end{align}
When we finally substitute this quantity back into the equation we can simplify it to
\begin{align}
	b_2(\Qf{0,2}+\Qf{0,1}) = (b_2(f_2^2-f_2^0)-a_2(1+t))\Qf{0,0}-(1+t)\Qc{0,0}\,.
\end{align}
The question is can we find the quantity $f_2^2-f_2^0$? It appears that
\begin{align}
	-(f_2^2-f_2^0) =1+a_1t+\frac{a_1^2}{b_1}(a_1^2-a_1-4)t^3 +\frac{4a_1^2}{b_1^2}(a_1^3-3a_1^2+1)t^4+\cdots
\end{align}
which is not a nice form. Unfortunately we haven't been able to go further with this case. One could try making some assumptions for the form of equations like these for $\Qf{m,n}$,$\Qc{m,n}$ and $\Qh{m,n}$ but even then it's not clear if one find a relation between $\Qf{0,0}$, $\Qc{0,0}$,$\Qh{0,0}$ and $\Qch{0,0}$. We would have 6 equations like the one above plus 6 more where $b_i = 0$ for the $b_i$ not present. Then one may use the 4 regular coefficient extractions from the master functional equation for $\Qf{m,n}$ and $m,n\leq 1$ plus 12 more where we set one or more of $b_1,b_2$ to zero. This brings us to a grand total of 28 equations in 36 unknown variables of which we need to eliminate 32. It's not entirely clear if we may do that.

\endinput
